\section{The PMNS Matrix: Neutrino Mixing Protocol}

\subsection{A. The Fundamental Difference: Knots vs. Phonons}

The most striking feature of the flavor sector is the magnitude disparity between quark mixing (CKM) and neutrino mixing (PMNS). As established in \textbf{Paper II}, this is a consequence of the structural duality of the substrate:

\begin{enumerate}
    \item \textbf{Quarks are Topological Solitons ($\Delta^n$):} Constrained by Manifold Drag.
    \item \textbf{Neutrinos are Lattice Phonons ($\Delta^{-n}$):} Characterized by \textbf{Symmetry Transparency}.
\end{enumerate}

Because neutrinos are delocalized vibrations ($M \approx 9$ meV), they do not interact with the "interface crack" ($\pi$) like quarks; they resonate across the global symmetry ratios of the lattice ($\sigma/\nu$). This explains why PMNS angles are large and geometric ($30^\circ, 45^\circ$) while CKM angles are small and frictional ($13^\circ$).

\subsection{B. The Solar Angle ($\theta_{12}$): The Symmetry Ratio}

\textbf{The Problem:} How does a delocalized phonon distribute itself across the lattice substrate?

\textbf{The Geometric Derivation:} Unlike fermions which are restricted to specific integer channels, phonons resonate across the full Interaction Order ($\sigma$) relative to the Chiral Basis ($\nu$). This creates a distribution governed by the projection geometry of the lattice.

\subsubsection{1. First Order Approximation (Integer Hardware)}
At the bit-level of the lattice hardware, the mixing is defined by the integer ratio of the Symmetry Order to the Channel Capacity.

\begin{itemize}
    \item \textbf{Numerator ($\sigma = 5$):} The 5 interaction shells of the Petrie projection.
    \item \textbf{Denominator ($\nu = 16$):} The fundamental bit-depth of the chiral capacity.
\end{itemize}

\begin{equation}
\sin^2 \theta_{12} \approx \frac{\sigma}{\nu} = \frac{5}{16} = \mathbf{0.3125}
\end{equation}
\textit{(Error: 1.8\% relative to experiment)}

\subsubsection{2. The Phonon Refinement (Golden Geometry)}
While the integer ratio provides a strong first approximation, neutrinos are lattice phonons (waves), not static knots. They couple to the \textit{continuous geometry} of the 4D projection rather than the discrete bits.

The projection of $E_8$ to 4D is governed by the $H_4$ Coxeter group (the 600-cell), where the 5-fold symmetry manifests continuously as the Golden Ratio ($\phi$). The exact geometric solution is the inverse of the Golden aperture:

\begin{equation}
\sin^2 \theta_{12} = \frac{1}{2\phi} \approx \frac{1}{3.236} = \mathbf{0.3090}
\end{equation}

\begin{itemize}
    \item \textbf{Experimental Target (NuFIT 5.3):} $0.307 \pm 0.013$
    \item \textbf{Accuracy:} \textbf{99.4\%} (0.6\% deviation)
\end{itemize}

\textbf{Physical Meaning:} The Solar angle is the ratio of Lattice Symmetry to Capacity. Neutrinos, being waves, "see" the continuous Golden Geometry ($1/2\phi$) of the substrate, while Quarks "see" the discrete Integer Channels ($5/16$).

\subsection{C. The Atmospheric Angle ($\theta_{23}$): The Boundary Limit}

\textit{The Problem:} What is the maximum mixing entropy allowed by the topology?

\textit{The Geometric Derivation:} The mixing between the second and third generations represents the maximal delocalization allowed by the Topological Boundary ($\chi$).

\begin{itemize}
    \item \textbf{Constraint ($\chi = 2$):} From the Mapping Function, $\chi$ is the Euler characteristic defining the node.
    \item \textbf{The Limit:} In a system with a 2-point boundary, the maximum Shannon entropy state is an equal partition ($1/2$).
\end{itemize}

\begin{equation}
\sin^2 \theta_{23} = \frac{1}{\chi} = \frac{1}{2} = \mathbf{0.5000}
\end{equation}

\begin{itemize}
    \item \textbf{Experimental Target (NuFIT 5.3):} $0.51 \pm 0.03$ (Normal Hierarchy)
    \item \textbf{Accuracy:} 98.0\%
    \item \textbf{Physical Meaning:} This is Maximal Mixing. The $\chi=2$ boundary allows equal probability distribution between the $\mu$ and $\tau$ sectors. This explains why $\theta_{23}$ is close to $45^\circ$ while quark angles are small; quarks are constrained by the Manifold Drag ($\pi$), while neutrinos are constrained only by the Topology ($\chi$).
\end{itemize}

\subsection{D. The Reactor Angle ($\theta_{13}$): The Interaction Leak}

\textit{The Problem:} Why is $\theta_{13}$ non-zero but small?

\textit{The Geometric Derivation:} This angle represents the leakage of the Interaction Remainder (the quark/color sector) into the neutrino sector via the vacuum impedance.

\begin{itemize}
    \item \textbf{Numerator ($\sigma - \chi = 3$):} The \textbf{Interaction Remainder} derived in Paper I. This is the specific structural invariant that mandates exactly three quark colors and three fermion generations.
    \item \textbf{Denominator ($\alpha^{-1} \approx 137$):} The Vacuum Impedance.
\end{itemize}

\begin{equation}
\sin^2 \theta_{13} = \frac{\sigma - \chi}{\alpha^{-1}} = \frac{3}{137.036} \approx \mathbf{0.02189}
\end{equation}

\begin{itemize}
    \item \textbf{Experimental Target (NuFIT 5.3):} $0.0220 \pm 0.0006$
    \item \textbf{Accuracy:} 99.5\%
    \item \textbf{Physical Meaning:} The Reactor angle is the Sector Interference between the Matter Sector (3) and the Vacuum Field ($\alpha^{-1}$). It proves that the "3" appearing in the PMNS matrix is the same "3" that defines the number of generations and colors.
\end{itemize}

\subsection{E. Geometric Predictions}

The identification of neutrinos as Lattice Phonons rather than topological knots imposes strict constraints on their physical behavior, distinct from standard field theory expectations.

\begin{enumerate}
    \item \textbf{Neutrino Mass Hierarchy (Normal Ordering):}
    \begin{itemize}
        \item \textbf{Prediction:} The JUNO (2026) and DUNE experiments will confirm the Normal Hierarchy ($m_1 < m_2 < m_3$).
        \item \textbf{Mechanism:} The system fills volumetric resonances from the ground state upward. An Inverted Hierarchy ($m_3 < m_1$) would imply a "Topological Debt" where higher-order lattice excitations exist without the foundational support of the lower tiers, violating the Lagrangian of Persistence.
    \end{itemize}

    \item \textbf{The Cosmological Mass Discrepancy:}
    \begin{itemize}
        \item \textbf{Prediction:} While oscillation experiments measure the Protocol Latency (friction), the Euclid Satellite (2026) will measure the Structural Rest Mass, settling on a sum significantly below the current $60 \text{ meV}$ floor (predicted $\Sigma m_\nu \approx 0.20 \text{ meV}$).
        \item \textbf{Mechanism:} Because phonons decouple from the Manifold Drag ($\pi D$), their gravitational mass (rest mass) is orders of magnitude lower than their inertial resistance to flavor change (oscillation mass). This discrepancy is the signature of a Phonon-based neutrino sector.
    \end{itemize}
\end{enumerate}